% Pandoc template for the fed_fraud documents.
% Based on proposal_gen/templates/detailed_proposal_template/detailed_1_classic.tex
% (Classic report): centered title page, table of contents, ruled section
% headings and the Geodesic house footer. Fonts are Latin Modern (lmodern).
% Two layouts: the full report (title page and contents) and, with
% -V brief, a compact one-page-brief header in the same style.
\documentclass[letterpaper,11pt]{article}
\usepackage[T1]{fontenc}
\usepackage[utf8]{inputenc}
\usepackage{lmodern}
\usepackage[letterpaper,top=0.9in,bottom=1.2in,left=1in,right=1in,footskip=0.6in]{geometry}
\usepackage[table]{xcolor}
\usepackage{geodesicbrand}
\usepackage{amsmath,amssymb}
\usepackage{enumitem}
\usepackage{array}
\usepackage{calc}
\usepackage{longtable}
\usepackage{makecell}
\usepackage[most]{tcolorbox}
\usepackage{titlesec}
\usepackage{needspace}
\usepackage{etoolbox}
\usepackage{hyperref}
\hypersetup{colorlinks=true,urlcolor=coolblack,linkcolor=coolblack,filecolor=geoblue,
  pdftitle={Federal fraud expansion},pdfauthor={\geofirm}}
\urlstyle{same}

\setlist{itemsep=2pt,leftmargin=1.4em}
\setlength{\parindent}{0pt}
\setlength{\parskip}{6pt}
\providecommand{\tightlist}{\setlength{\itemsep}{0pt}\setlength{\parskip}{0pt}}
\setcounter{secnumdepth}{2}

% Section style: house rule under each heading
\titleformat{\section}{\Large\bfseries\color{coolblack}}{\thesection}{0.6em}{}[{\color{coolblack}\titlerule[0.5pt]}]
\titleformat{\subsection}{\large\bfseries\color{geomid}}{\thesubsection}{0.6em}{}
\titleformat{\subsubsection}{\normalsize\bfseries\color{geomid}}{\thesubsubsection}{0.6em}{}
\titlespacing*{\section}{0pt}{16pt}{8pt}

% Tables are written by latex/filters.lua (navy header, banded rows)
\setlength{\LTpre}{6pt}\setlength{\LTpost}{6pt}

\begin{document}

% TITLE PAGE
\begin{titlepage}
\newcommand{\HRule}{\rule{\linewidth}{0.3mm}}
\begin{center}
\vspace*{0.2in}
\textsc{\LARGE \geofirm}\\[0.6cm]
\includegraphics[width=30mm]{icon.png}\\[8mm]
\textsc{\Large Federal Fraud Expansion\\[0.15cm] Working Document}\\[0.5cm]
\HRule\\[0.4cm]
{\huge\bfseries Federal fraud expansion}\\[0.3cm]
{\large Working plan}\\[0.2cm]
\HRule\\[1cm]
\begin{minipage}{0.8\textwidth}\centering\large
Authored by:\\[0.18cm]
\geosigner, \geosignertitle: \geoemail
\end{minipage}\\[0.8cm]
{\large When: \compiletime}\\[1.2cm]
\geosurface{catenoid}{4.6cm}\\[0.9cm]
$\mathcal{G}(x)=x^x\qquad\nabla_{\dot\gamma}\dot\gamma=0$\\[3mm]
{\small\itshape \geopromise}
\end{center}
\end{titlepage}

\pagestyle{geodesic}
\tableofcontents
\newpage

\textbf{Status:} Draft v0.1 (2026-10-25). Waiting on Steve's edits
before step 1 starts.

\hypertarget{goal}{%
\section{Goal}\label{goal}}

Take the SBA fraud work done with Rosemary and turn it into a repeatable
playbook. Apply it first to the agencies in the federal fraud initiative
(america.gov/fraud), then use those pilots to make the case for a
government-wide program.

\hypertarget{starting-assumptions-please-correct}{%
\section{Starting assumptions (please
correct)}\label{starting-assumptions-please-correct}}

\begin{itemize}
\tightlist
\item
  \textbf{Starting agencies.} We could not read america.gov/fraud
  directly because the site blocks automated access. This list comes
  from public coverage of the March 2026 anti-fraud task force executive
  order and DOJ's National Fraud Detection Center. It needs to be
  checked against the site.

  \begin{itemize}
  \tightlist
  \item
    Agencies named under the task force: DOJ, Treasury, HHS, Education,
    HUD, Labor, DHS, SBA and OMB, plus the Inspectors General.
  \item
    Inspector General offices in the Detection Center: USDA, Education,
    HHS, DHS, HUD, Interior, Labor, VA, DoD (DCIS), TIGTA, SBA and SSA.
  \end{itemize}
\item
  \textbf{Data.} We start with public data: USAspending, SAM.gov, FPDS,
  PRAC data, OIG/GAO reports and DOJ press releases. Agencies would
  provide internal data only after a pilot is approved.
\item
  \textbf{Audience.} Proposals go to agency leadership, such as deputy
  secretaries, CFOs and chief data officers, and their Inspectors
  General. They also go to the task force or OMB as the cross-agency
  sponsor.
\item
  \textbf{What the SBA work produced.} Not yet in this repo. Step 1
  depends on getting those materials.
\end{itemize}

\hypertarget{plan}{%
\section{Plan}\label{plan}}

\hypertarget{step-1-turn-the-sba-work-into-a-reusable-playbook}{%
\subsection{Step 1: Turn the SBA work into a reusable
playbook}\label{step-1-turn-the-sba-work-into-a-reusable-playbook}}

Write down what the SBA effort does. That covers its data sources, fraud
typologies (for example identity theft, shell entities, duplicate or
stacked awards, and ineligible recipients), detection methods, outputs
and measured results. Then separate the parts that carry over to other
agencies from the parts that only fit SBA.

\textbf{Deliverable:} \texttt{docs/playbook/sba-baseline.md} and a
one-page summary of the method.

\hypertarget{step-2-profile-the-initiative-agencies-and-rank-them}{%
\subsection{Step 2: Profile the initiative agencies and rank
them}\label{step-2-profile-the-initiative-agencies-and-rank-them}}

For each agency, note its largest programs and how much money they pay
out. Also note its known fraud and improper-payment exposure, using
PaymentAccuracy.gov and OIG/GAO findings. Record which public data
exists for it and how closely it matches the SBA playbook. Rank the
agencies by impact and feasibility, and choose 2 or 3 for the first
wave.

\textbf{Deliverable:} \texttt{docs/agencies/landscape.md}, with one
short profile per agency and a ranking matrix.

\hypertarget{step-3-prototype-against-public-data-for-the-first-wave-agencies}{%
\subsection{Step 3: Prototype against public data for the first-wave
agencies}\label{step-3-prototype-against-public-data-for-the-first-wave-agencies}}

Apply the SBA detection patterns to public data for the chosen agencies.
Build wireframes for the analyst and leadership views. The goal is a
convincing demo built from real public data, not a production system.

\textbf{Deliverable:} code under
\texttt{prototypes/\textless{}agency\textgreater{}/} and wireframes
under \texttt{docs/wireframes/}.

\hypertarget{step-4-write-proposals-for-agency-leadership}{%
\subsection{Step 4: Write proposals for agency
leadership}\label{step-4-write-proposals-for-agency-leadership}}

Write one proposal per first-wave agency. Each covers the problem and
estimated dollar exposure, what the prototype found, the pilot scope and
data request, timeline, cost, success metrics and governance. Add a
cross-agency proposal that presents the shared platform as the route to
government-wide coverage.

\textbf{Deliverable:}
\texttt{docs/proposals/\textless{}agency\textgreater{}.md} and a slide
deck for each meeting.

\hypertarget{open-questions-for-steve}{%
\section{Open questions for Steve}\label{open-questions-for-steve}}

\begin{enumerate}
\def\labelenumi{\arabic{enumi}.}
\tightlist
\item
  What does the SBA work produce today, and can it be added here?
\item
  Who is the first audience: a specific agency leader, an IG, OMB, or
  the task force?
\item
  Is any non-public data access expected for the pilots?
\item
  Is there a deadline or meeting the first output has to be ready for?
\end{enumerate}

\hypertarget{sources}{%
\section{Sources}\label{sources}}

\begin{itemize}
\tightlist
\item
  \href{https://www.bakermckenzie.com/en/insight/publications/2026/03/united-states-trump-launches-nationwide-antifraud-task-force}{Baker
  McKenzie: Trump launches nationwide antifraud task force (Mar 2026)}
\item
  \href{https://www.executivegov.com/articles/doj-national-fraud-detection-center}{ExecutiveGov:
  DOJ launches National Fraud Detection Center}
\end{itemize}

\end{document}
